% Part: sets-functions-relations
% Chapter: infinite
% Section: dedekinds-proof
%
\documentclass[../../../include/open-logic-section]{subfiles}

\begin{document}

\olfileid{sfr}{infinite}{dedekindsproof}

\olsection[ڈیڈیکائنڈ دا ”ثبوت“]{لامتناہی سیٹ دے وجود دا ڈیڈیکائنڈ دا
”ثبوت“}

ایس باب وچ اسی ڈیڈیکائنڈ الجبریاں دے حوالے نال قدرتی عدداں دا
سیٹ-تھیوری بیان دِتا اے۔ \olref[arith][ref]{sec} وچ اسی (ہور گلاں دے
نال) تجزیے دی اعدادی تشکیل دی فلسفیانہ اہمیت اُتے غور کیتا سی۔ ہُن
سانوں ایتھے حاصل کیتی کامیابی دی اہمیت اُتے غور کرنا چاہیدا اے۔

پورے \olref[sfr][arith][]{chap} وچ اسی قدرتی عدداں نوں پہلے توں ملیا
منیا، تے اوہناں نال صحیح، ناطق تے حقیقی عدد صاف طور تے بنائے۔ ایس باب
وچ اسی قدرتی عدداں دی کوئی صاف بناوٹ نہیں دِتی۔ اسی صرف ایہ وکھایا اے
کہ \emph{کوئی وی ڈیڈیکائنڈ لامتناہی سیٹ دِتا ہووے}، تاں اسی اک ایسا
سیٹ تعریف کر سکدے آں جیہڑا عین اوویں ورتاؤ کرے گا جیویں اسی چاندے آں
کہ~$\Nat$ ورتاؤ کرے۔

سو صاف اے کہ اسی کسے مابعدالطبیعیاتی سوال، مثلاً
\emph{قدرتی عدد کیہڑیاں شےواں نیں}، دا جواب دین دا دعویٰ نہیں کر سکدے۔
پر ایہ چنگی گل اے۔ آخر \olref[sfr][arith][ref]{sec} وچ اسی زور دِتا سی
کہ $\Real$ نوں ڈیڈیکائنڈ کٹاں دا سیٹ تعریف کرن نوں اک سہولت والی
مقرری دی بجائے \emph{دریافت} سمجھنا غلط ہووے گا۔ فیصلہ کن مشاہدہ ایہ اے
کہ ڈیڈیکائنڈ کٹ حقیقی عدداں دیاں اہم ریاضیاتی خاصیتاں دی مثال بن دے
نیں۔ ایتھے وی ایسے طرح: فیصلہ کن مشاہدہ ایہ اے کہ \emph{ہر} ڈیڈیکائنڈ
الجبرا قدرتی عدداں دیاں اہم ریاضیاتی خاصیتاں دی مثال بن دا اے۔ (اصل وچ
ڈیڈیکائنڈ نے ایہ ثابت کر کے ایس نکتے نوں پکا کیتا کہ سارے ڈیڈیکائنڈ
الجبرے \emph{ہم شکل} نیں (\citeyear[Theorems
132--3]{Dedekind1888})۔ سو ایہ حیرانی دی گل نہیں کہ کئی اجوکے
”ساختیت پسند“ ڈیڈیکائنڈ نوں اپنا پیش رو آکھدے نیں۔)

نال ای، اسی وکھایا اے کہ قدرتی عدداں دے نظریے نوں اک سادہ، ناپختی سیٹ
تھیوری وچ کیویں سمایا جا سکدا اے؛ ایہ سیٹ تھیوری آپ ہجے خاصی غیر رسمی
اے، پر ایہ (بظاہر) قدرتی عدداں نوں پہلے توں ملیا نہیں مندی۔ ایس لئی
شاید اسی ڈیڈیکائنڈ دے اپنے ولولے بھرے منصوبے نوں پورا کرن دے راہ اُتے
آں، جیہنوں اوہنے ایویں بیان کیتا:
\begin{quote}
	سائنس وچ جیہڑی گل ثبوت دے قابل اے، اوہنوں ثبوت توں بناں منّنا
	نہیں چاہیدا۔ بھاویں ایہ تقاضا معقول لگدا اے، میں ایہ نہیں من سکدا
	کہ سادہ ترین سائنس دیاں بنیاداں رکھن دے سب توں نویں طریقیاں وچ وی
	ایہ پورا ہویا اے؛ یعنی منطق دا اوہ حصہ جیہڑا عدداں دے نظریے نال
	واسطہ رکھدا اے۔ حساب (الجبرا، تجزیہ) نوں محض منطق دا اک حصہ آکھ کے
	میرا مطلب ایہ اے کہ میں عدد دے تصور نوں تھاں تے سمے دے تصورات یا
	وجدان توں پوری طرح آزاد سمجھدا آں---بلکہ میں اوہنوں فکر دے خالص
	قانوناں دی فوری پیداوار سمجھدا آں۔
	\citep[preface]{Dedekind1888}
\end{quote}
ڈیڈیکائنڈ دا دلیر خیال ایہ اے۔ اسی ہُن (سادی) سیٹ تھیوری اکیلی نال قدرتی
عدد بناؤن دا طریقہ وکھایا اے۔ \olref[sfr][arith][]{chap} وچ اسی ویکھیا
سی کہ قدرتی عدد تے کجھ سیٹ تھیوری دِتے ہوون، تاں حقیقی عدد کیویں بنائے
جا سکدے نیں۔ سو شاید ”حساب (الجبرا، تجزیہ)“، ”محض منطق دا اک حصہ“
نکلے (لفظ ”منطق“ دے ڈیڈیکائنڈ والے وسیع معنٰی وچ)۔

خیال ایہو اے۔ پر اک پل رکو۔ ڈیڈیکائنڈ الجبرے (قدرتی عدداں دے ساڈے
قائم مقام) دی ساڈی بناوٹ ڈیڈیکائنڈ لامتناہی سیٹ دے وجود اُتے مشروط اے۔
(ذرا پچھے \olref[sfr][infinite][dedekind]{thm:DedekindInfiniteAlgebra}
ول ویکھو۔) جے ڈیڈیکائنڈ لامتناہی سیٹ دا وجود ”منطق“ یا ”فکر دے خالص
قانوناں“ راہیں قائم نہ کیتا جا سکے، تاں منصوبہ ایتھے ای رک جاندا اے۔

سو، کی ڈیڈیکائنڈ لامتناہی سیٹ دا وجود ”فکر دے خالص قانوناں“ نال
\emph{قائم} کیتا جا سکدا اے؟ ڈیڈیکائنڈ دی کوشش ایہ سی:
\begin{quote}
  میرے خیالاں دی اپنی قلمرو، یعنی اوہناں ساریاں شےواں دی کُلّیت $S$
  جیہڑیاں میرے خیال دا موضوع بن سکدیاں نیں، لامتناہی اے۔ کیوں جو جے
  $s$،~$S$ دے اک رکن نوں ظاہر کردا اے، تاں ایہ خیال $s'$ کہ~$s$ میرے
  خیال دا موضوع بن سکدا اے، آپ~$S$ دا اک رکن اے۔ جے اسی ایہنوں
  عکس $\phi(s)$، یعنی رکن~$s$ دا عکس، سمجھئیے، تاں \dots~$S$ [ڈیڈیکائنڈ]
  لامتناہی اے، تے ایہی ثابت کرنا سی۔
	\citep[\S66]{Dedekind1888}
\end{quote}
اک ایسی کتاب دے عین وچکار، جیہڑی زیادہ تر نہایت کڑے ریاضیاتی ثبوتاں
اُتے مشتمل اے، ایہ گل لبھنا خاصا حیران کن اے۔ دو نکات قابلِ ذکر نیں۔

پہلا: ایس ”ثبوت“ وچ شاید ہی اوہ کردار اے جیہنوں اسی اج ”ریاضیاتی“
منّاں گے۔ ایہ نفسیاتی شےواں (خیالاں) دی گل کردا اے، تے اوہ وی محض
\emph{ممکن} خیالاں دی۔

دوجا: گھٹ توں گھٹ ڈیڈیکائنڈ الجبرے دی ساڈی پیشکش دے مطابق، ایس ”ثبوت“
وچ اک صاف فنی کمی اے۔ جے ڈیڈیکائنڈ دی دلیل کامیاب اے، تاں ایہ صرف
اینا قائم کردی اے کہ شےواں دی تعداد لامتناہی اے (خاص طور تے خیالاں دی
تعداد لامتناہی اے)۔ پر ڈیڈیکائنڈ نوں سانوں ایہ وجہ وی دینی
پینی اے کہ~$S$ نوں اک اکیلا \emph{سیٹ} سمجھیا جاوے، جیہدے لامتناہی
!!{element}s نیں، نہ کہ~$S$ نوں \emph{کجھ شےواں} سمجھیا جاوے (جمع
وچ)۔

ڈیڈیکائنڈ نے ایتھے کوئی خلا نہ ویکھیا، ایس حقیقت توں شاید ایہ پتا لگے
کہ ”کُلّیت“ لفظ دی اوہدی ورتوں عین \emph{ساڈی} ”سیٹ“ لفظ دی ورتوں
نال نہیں رلدی۔\footnote{اصل وچ ساڈے کول ایہ سوچن دیاں ہور وجوہ وی نیں
کہ ایہ نہیں رلدی سی؛ \citet[p.~23]{Potter2004} ویکھو۔} پر ایہ کوئی
وڈی حیرانی نہیں ہووے گی۔ پچھلے دو باباں وچ ساڈا منصوبہ---قدرتی عدداں
دی اک ”بناوٹ“، تے اوہناں توں صحیح، حقیقی تے ناطق عدداں دی اک
”بناوٹ“---سارا سادہ طور تے چلایا گیا اے۔ جدوں جدوں سانوں لوڑ پئی، اسی
ایہ سیٹ یا اوہ سیٹ اپنے کم لئی لے لیا، تے ایہ عام اصول بہت گھٹ دِتے کہ
عین کیہڑے سیٹ کدوں موجود ہوندے نیں۔ پر اسی جان دے آں کہ سانوں
\emph{کجھ} عام اصول چاہیدے نیں، نہیں تاں اسی رَسل دے تضاد وچ جا
ڈِگاں گے۔

ہُن ساڈی سادہ لوحی توں اگے ودھن دا ویلا آ گیا اے۔

\end{document}
